% Preamble bersama untuk kedua laporan TK 2 (soal1 dan soal2).
% Mengikuti "Petunjuk Penulisan Dokumen Project Anum":
%   - A4, font standar LaTeX (Latin Modern), 12pt, margin >= 2.54 cm
%   - judul, rangkuman, daftar isi TIDAK dinomori; bagian lain dinomori dari 1
%   - tabel/gambar/kode/persamaan dinomori mengikuti bagian: Tabel 3.1, Gambar 3.2, Kode 3.1
\usepackage[T1]{fontenc}
\usepackage[utf8]{inputenc}
\usepackage[indonesian]{babel}
\usepackage{lmodern}
\usepackage[a4paper,margin=2.54cm]{geometry}
\usepackage{amsmath,amssymb,booktabs,array,graphicx,float,microtype}
\usepackage{algorithm,algpseudocode}
\usepackage{listings}
\usepackage[hidelinks]{hyperref}
\emergencystretch=1em

\numberwithin{equation}{section}
\counterwithin{table}{section}
\counterwithin{figure}{section}
\counterwithin{algorithm}{section}
\floatname{algorithm}{Kode}
\renewcommand{\lstlistingname}{Kode}
\algrenewcommand\algorithmicrequire{\textbf{Masukan:}}
\algrenewcommand\algorithmicensure{\textbf{Keluaran:}}
\setlength{\belowcaptionskip}{6pt}
\lstset{language=Python, basicstyle=\ttfamily\footnotesize, breaklines=true,
  columns=fullflexible, keepspaces=true, showstringspaces=false, frame=single}

\newcommand{\R}{\mathbb{R}}
\newcommand{\T}{\mathsf{T}}
\newcommand{\norm}[1]{\left\lVert #1\right\rVert_2}
\newcommand{\bx}{\mathbf{x}}
\newcommand{\bd}{\mathbf{d}}
\newcommand{\bg}{\mathbf{g}}
\newcommand{\bs}{\mathbf{s}}
\newcommand{\by}{\mathbf{y}}

% \input biasa di dalam tabular gagal bila diikuti \bottomrule; pakai primitif \@@input
\makeatletter\newcommand{\isitabel}[1]{\@@input #1 }\makeatother

% Grafik dicari di report/figures/ dan results/ (relatif terhadap folder report/)
\graphicspath{{figures/}{../results/}}
